\section{Quadratic orders and infinitely many torsion classes}
\label{sec:singular-orders}

Failure of generation can already occur over a one-dimensional local
domain whose generic fibre is a field.  The obstruction in this section
is a family of additive invariants of integral involutions.  It records
information that disappears from every ordinary character.

\begin{theorem}\label{thm:singular-order}
Let $O=\Z_{(2)}$ or $\Z_2$, let $r,s\in O$, and suppose
$t^2+2st+4r$ is irreducible over $\operatorname{Frac}(O)$.  Put
\[
 A=O[t]/(t^2+2st+4r).
\]
Then $A$ is a one-dimensional noetherian local characteristic-zero
$\Z$-flat domain, and $\Phi_A$ is not surjective.  More precisely, the
kernel of
\[
 K_0(\Lat_A(C_2))\longrightarrow
 K_0(\operatorname{Rep}_{\operatorname{Frac}(A)}(C_2))
\]
contains a direct summand isomorphic to a countable direct sum of
copies of $\Z/2$.
\end{theorem}

For $r=1$ and $s=0$, one nonzero class of exact order two is
\[
 [L]-[A[C_2]],\qquad L=A^2,\qquad
 T_L=\begin{pmatrix}1+t&2\\2-t&-1-t\end{pmatrix}.
\]


We first construct the invariants, including their exact-sequence
compatibility.  All reductions in the proof are modulo $2$ in $O$;
write $k=\F_2$.

\begin{lemma}\label{lem:order-detectors}
Let $\mathcal F_s$ be the monic irreducible polynomials in $k[z]$
other than $z$, also excluding $z+1$ when $s$ is odd.  For every
$f\in\mathcal F_s$, there is an
additive homomorphism
\[
 \nu_f:K_0(\Lat_A(C_2))\longrightarrow\Z/2
\]
that vanishes on every lattice obtained by extending scalars from $O$.
\end{lemma}

\begin{proof}
Put $n=\operatorname{rank}_A E$ and choose an $A$-basis.  Its involution has a unique
expression $T=P+tQ$ with $P,Q$ matrices over $O$.  Since $1,t$ is an
$O$-basis of $A$, the equation $T^2=I$ gives
\[
 P^2-4rQ^2=I,\qquad PQ+QP=2sQ^2.
\]
Set $D=P-I$.  Then
\begin{equation}\label{eq:order-curvature}
 D^2+2D=4rQ^2,\qquad DQ+QD=2sQ^2-2Q.
\end{equation}
The reduction $d$ of $D$ has square zero on $V=k^n$.
Put $H(E)=\ker d/\operatorname{im}d$.  The reduction of $Q$ induces an
operator $B$ on $H(E)$.  For a class represented by $v\in O^n$ with
$Dv\in2O^n$, define
\[
 C[v]=[Dv/2\bmod2].
\]
This is well-defined.  The vector $Dv/2$ is a cycle because
$D(Dv/2)=2rQ^2v-Dv$ is even.  Changing $v$ by $2w$ changes its image by
the boundary $Dw$; changing its homology class by the boundary $Dw$
changes the image by $D^2w/2=2rQ^2w-Dw$, again a boundary modulo two.
Equation~\eqref{eq:order-curvature} now gives
\begin{equation}\label{eq:order-operators}
 C^2+C=\overline r B^2,\qquad CB+BC=B+\overline s B^2.
\end{equation}
In particular $Z=B^2$ commutes with both operators.

These operators are natural.  If an equivariant map is written
$F=F_0+tF_1$, comparison of coefficients gives
\[
 P'F_0-F_0P=4r(Q'F_1-F_1Q),
 \]
 \[
 Q'F_0-F_0Q=F_1P-P'F_1+2s(Q'F_1-F_1Q).
\]
Thus $\overline F_0$ induces a map on $H$ commuting with $B$; the second
identity makes its commutator a boundary on cycles.  The first identity,
after division by two on cycles, proves commutation with $C$.
In particular the construction is independent of the chosen basis.

Fix $f\in\mathcal F_s$ of degree $a$, and let $H_f(E)$ be the
$f$-primary subspace for $Z$.  Its dimension is divisible by $2a$.
Indeed, filtration by powers of $f(Z)$ reduces the assertion to
modules over $F=k[z]/(f)$ on which $Z$ is $\alpha=z\bmod f$.
Set $U=B+\overline s\alpha$.  Equations~\eqref{eq:order-operators}
give
\[
 CU+UC=U,\qquad U^2=\eta:=\alpha+\overline s\alpha^2,
 \qquad C^2+C=\overline r\alpha.
\]
The exclusions in $\mathcal F_s$ ensure $\eta\ne0$.
This operator algebra is $\operatorname{Mat}_2(F)$: with
$\beta^2=\eta$ in the finite field $F$, send
\[
 U/\beta\longmapsto\begin{pmatrix}1&0\\1&1\end{pmatrix},\qquad
 C\longmapsto\begin{pmatrix}0&1\\\overline r\alpha&1\end{pmatrix}.
\]
The defining relations hold, and the four spanning elements
$1,U,C,UC$ map to independent matrices.
Every finite module therefore
has $F$-dimension divisible by two.  Define
\begin{equation}\label{eq:order-invariant}
 \nu_f(E)=\frac{\dim_k H_f(E)}{2\deg f}\pmod2.
\end{equation}

It remains to prove additivity; dimension of homology alone would not
suffice.  In an $A$-split exact sequence
$0\to E'\to E\to E''\to0$, choose compatible bases and write
\[
 P=\begin{pmatrix}P'&U\\0&P''\end{pmatrix},\qquad
 Q=\begin{pmatrix}Q'&W\\0&Q''\end{pmatrix}.
\]
The square-zero reductions give the exact three-cycle
\begin{equation}\label{eq:order-three-cycle}
 H(E')\longrightarrow H(E)\longrightarrow H(E'')
       \xrightarrow{\partial}H(E').
\end{equation}
The first two maps commute with $B,C$.  The connecting map satisfies
\begin{equation}\label{eq:order-connecting}
 B'\partial=\partial B'',\qquad
 C'\partial=\partial(C''+I).
\end{equation}
For a cycle $v$ in the last term, $\partial[v]=[Uv]$.  The off-diagonal
parts of~\eqref{eq:order-curvature} are
\[
 D'U+UD''+2U=4r(Q'W+WQ''),
 \]
 \[
 Q'U+UQ''=-D'W-WD''-2W+2s(Q'W+WQ'').
\]
Division of the first identity by two on a cycle proves the formula
for $C$ in~\eqref{eq:order-connecting}; the second proves the formula
for $B$ modulo boundaries.

All three maps preserve the primary subspaces for $Z$, and all their
images are invariant under $B,C$.  Their image dimensions are therefore
multiples of $2a$, by the matrix-algebra calculation.  Primary
decomposition preserves exactness in~\eqref{eq:order-three-cycle}.
If its three image dimensions are $u,v,w$, then the sum of its three
term dimensions is $2(u+v+w)$, a multiple of $4a$.
Dividing by $2a$ modulo two gives
$\nu_f(E)=\nu_f(E')+\nu_f(E'')$.
Finally, for a lattice defined over $O$ one has $Q=0$, hence $B=0$ and
$H_f=0$ for every $f\ne z$.
\end{proof}

\begin{proof}[Proof of Theorem~\ref{thm:singular-order}]
For each $f\in\mathcal F_s$, let $a=\deg f$ and choose the
matrix $D_f$ with entries $0,1$ lifting the companion matrix of $f$.  Define
\[
 B_f=\begin{pmatrix}D_f&0\\sD_f^2-D_f&-D_f\end{pmatrix},\qquad
 C_f=\begin{pmatrix}0&I\\rD_f^2&-I\end{pmatrix},\qquad
 T_f=I+2C_f+tB_f.
\]
Direct multiplication gives
$C_f^2+C_f=rB_f^2$ and $B_fC_f+C_fB_f=sB_f^2-B_f$, so that $T_f^2=I$ over
$A$.  Let $L_f=A^{2a}$ with this action.  Its square-zero reduction is
zero, and on $H(L_f)=k^{2a}$ the central operator is
\[
 Z=\operatorname{diag}(\overline D_f^{\,2},\overline D_f^{\,2}).
\]
Squaring is an automorphism of the field $k[z]/(f)$, so
$\overline D_f^{\,2}$ still has irreducible polynomial $f$.  Consequently
\[
 \nu_g(L_f)=\begin{cases}1,&g=f,\\0,&g\ne f.\end{cases}
\]

For every lattice $E$ there is an exact sequence
\[
 0\longrightarrow E\otimes\mathrm{sgn}
 \longrightarrow A[C_2]\otimes_A E
 \longrightarrow E\longrightarrow0,
\]
with diagonal action in the middle.  This middle lattice is a free
$A[C_2]$-module of rank $\operatorname{rank}_A E$, by untwisting the
action on the regular factor.  The matrix
\[
 J_f=\begin{pmatrix}I&0\\-I-tD_f&-I\end{pmatrix}
\]
is invertible and satisfies $J_fT_f=-T_fJ_f$.  Hence
$L_f\simeq L_f\otimes\mathrm{sgn}$ and
\[
 \delta_f=[L_f]-a[A[C_2]],\qquad 2\delta_f=0.
\]
The invariants $\nu_g$ show that these classes are independent and
nonzero.  For each lattice only finitely many $\nu_f$ are nonzero, so
$(\nu_f)_f$ takes values in the direct sum and retracts the inclusion
of the classes $\delta_f$.  They therefore form a direct summand, also
inside the ordinary-character kernel.  There are infinitely many monic
irreducibles in $\mathcal F_s$.
Over the fraction field each $T_f$ has $a$ eigenvalues $1$ and $a$
eigenvalues $-1$: it squares to one and has trace zero.  Thus all
$\delta_f$ have zero ordinary character.  When $r=1$, $s=0$, and $f=z+1$, this is the displayed rank-two example.

In degree two the subring generated by the trivial classes is spanned
by $\rho_2$ and $\rho_1^2=[A[C_2]]$.  Both have action matrices over
$O$, and all $\nu_f$ vanish on them, whereas $\nu_f(L_f)=1$ for every $f\in\mathcal F_s$.
This proves the failure of surjectivity.

Finally, irreducibility over $\operatorname{Frac}(O)$ and the monic
presentation make $A$ a domain, finite free of rank two over $O$.
Its quotient by $2$ is $k[t]/(t^2)$, so it is local with maximal ideal
$\mathfrak m=(2,t)$ and dimension one.  The defining equation has no
linear term in the cotangent space: its embedding dimension is two,
so the ring is singular.  Freeness over $O$ gives $\Z$-flatness.
\end{proof}

\begin{corollary}\label{cor:quadratic-order-classification}
Let $K/\Q_2$ be a quadratic field extension, and let $A$ be a
$\Z_2$-order in $K$.  Then $\Phi_A$ is an isomorphism if and only if
$A$ is the maximal order.  For every nonmaximal order, the
ordinary-character kernel in degree two has a direct summand
isomorphic to a countable direct sum of copies of $\Z/2$.
\end{corollary}

\begin{proof}
Write the maximal order as $O_K=\Z_2\oplus\Z_2\omega$.
Such a basis exists because $1$ is primitive in the free rank-two
$\Z_2$-module $O_K$.  The finite-index submodule
$A/\Z_2\subseteq O_K/\Z_2$ is $2^n\Z_2\omega$ for a unique $n\ge0$.
Consequently $A=\Z_2+2^nO_K$.  If $n=0$, it is a discrete valuation
ring and the positive theorem applies.  If $n\ge1$, the generator
$t=2^n\omega$ has polynomial
\[
 t^2-2^n\operatorname{Tr}_{K/\Q_2}(\omega)t
       +2^{2n}\operatorname{N}_{K/\Q_2}(\omega).
\]
The linear coefficient is divisible by two and the constant
coefficient by four.  The polynomial is irreducible over $\Q_2$,
so Theorem~\ref{thm:singular-order} applies.
\end{proof}
